The data the we used for train our neural models come from simulation and real world experiment. 
In the general setting, we can describe our dataset as tuple in the form of $\left\{[x_{t}, u_{t}], [x_{t+1}]\right\}$, 
were $x_{t}$ is the state of the system at time $t$, $u_{t}$ is the control input at time $t$ and $x_{t+1}$ is the state of the system at time $t+1$.
\\

For the baseline (\textit{simple-nueral-network}), no constraint on the dataset form are required, while on the physics informed models specific requirement on the
input form are needed. For be able to apply the Pinn framework with control, \cite{Antonelo_2024} shown that the input must 
assumed to be constant in a sub-interval $[0, T]$. 
This assumption bring us to generate several trajectorries with this specific constraint, and also limit the 
possibility of application on the data obtained in the real-world experiment. More details about this constraint are explored in the Pinn-section.
\\

All the sintetic data has been obtained by integrate the system dynamics with a Runge-kutta 4th order method, while the real-world data has been extrapolated 
using the ROS interface of the quadrotor.

